% ============================================================
%  2.1  Artificial Intelligence
% ============================================================
\section{Artificial Intelligence}
\label{sec:bg-ai}

Throughout human history the aspiration to replicate human thought through
technology has been a persistent theme; written testimony of this desire dates
as far back as the II century B.C.E.\ in the myth of the bronze guardian Talos
and in its desire for immortality~\cite{talos}.
Talos furnishes the mythological foundation for the Greek concept of the
\emph{automaton}---an entity that moves of its own accord---and has been
invoked to illustrate the ability to solve problems autonomously, not through
the formal specification of rules, but by learning from experience which
choice is the most appropriate.

This inductive process, which enables the extraction of patterns from data,
is known as \emph{machine learning}.
As an inductive method, it is heavily dependent on the quantity and quality of
the data from which it learns, and in particular on the
\emph{representational form} of that data.
Initially, the selection of representations---called \emph{features}---had to
be performed manually by domain experts; this limitation was
subsequently overcome by \emph{deep learning},
in which complex concepts are represented hierarchically in terms of simpler
ones.
When visualized, this hierarchy of concepts forms a deep computational graph,
leading to the term \textbf{Deep Learning} (DL).
The fundamental principle of DL is to learn an efficient \emph{latent
representation} of the input data, one designed to capture the essential
semantic information required for a given task and to transform the problem
into a space where a solution can be more easily discovered.


% ------------------------------------------------------------
\subsection{The Perceptron}
\label{sec:perceptron}
% ------------------------------------------------------------

The cornerstone of modern neural networks is the \emph{perceptron}, the first
tangible implementation of an artificial learning neuron.
The concept was developed throughout the 1950s and 1960s by Frank
Rosenblatt~\cite{perceptron}, inspired by earlier work from Warren McCulloch
and Walter Pitts.

A perceptron takes several binary inputs, $x_1, x_2, \dots$, and produces a
single binary output.
Rosenblatt proposed a simple rule to compute this output by introducing
\textit{weights}, $w_1, w_2, \dots$, which are real numbers expressing the
relative importance of each input.
The neuron's output, either~$0$ or~$1$, is determined by whether the weighted
sum of its inputs, $\sum_j w_j x_j$, falls below or exceeds some
\textit{threshold value}.
Like the weights, the threshold is a real-valued parameter of the neuron.
This decision rule can be expressed algebraically as:
\begin{equation}
\label{eq:perceptron-threshold}
\text{output} = 
\begin{cases} 
0 & \text{if } \sum_j w_j x_j \le \text{threshold}, \\ 
1 & \text{if } \sum_j w_j x_j > \text{threshold}. 
\end{cases}
\end{equation}

To simplify the notation, two standard substitutions are made.
First, the summation is rewritten as a dot product,
$\mathbf{w} \cdot \mathbf{x} \equiv \sum_j w_j x_j$, where $\mathbf{w}$
and~$\mathbf{x}$ are vectors of the weights and inputs, respectively.
Second, the threshold is moved to the other side of the inequality and
replaced by the neuron's \textit{bias}, defined as
$b \equiv -\text{threshold}$.
The bias can be thought of as a measure of how easily the perceptron can be
made to output a~$1$, or, in biological terms, to ``fire.''
Using the bias, the perceptron rule becomes:
\begin{equation}
\label{eq:perceptron-bias}
\text{output} = 
\begin{cases} 
0 & \text{if } \mathbf{w} \cdot \mathbf{x} + b \le 0, \\ 
1 & \text{if } \mathbf{w} \cdot \mathbf{x} + b > 0. 
\end{cases}
\end{equation}

This mathematical model can be viewed as a simple decision-making unit that
evaluates and combines multiple pieces of evidence through weighted
contributions.
By tuning its weights and bias, the perceptron can be adapted to represent
various types of decision processes.
Its ability to \emph{learn} from data by iteratively adjusting these
parameters constitutes the key advancement over the McCulloch--Pitts (MP)
neuron model, which relied solely on fixed, hand-designed weights and
thresholds.


% ------------------------------------------------------------
\subsection{Neural Networks and Activation Functions}
\label{sec:neural-networks}
% ------------------------------------------------------------

While a single perceptron can weigh evidence to make a simple (linear)
decision, its power is amplified when multiple perceptrons are organized into
a \emph{network} of layers.
In such a network, the first layer makes decisions based on the raw
input evidence.
The outputs of this layer are then fed as inputs to a second layer, which can,
in turn, make decisions at a more complex and abstract level.
By stacking multiple layers, the network can engage in highly sophisticated
decision-making.

From an abstract perspective, a neural network can be represented as a
directed acyclic graph composed of interconnected nodes---or
\emph{neurons}---inspired by their biological counterparts.
Each neuron receives a set of input signals, each associated with a
specific weight, and produces an output signal as a learned response.
The graph structure enables data propagation from input neurons to output
neurons through a series of weighted connections that modulate signal
strength at each stage.

\paragraph{Activation functions.}
Despite its foundational importance, the perceptron as presented in
Section~\ref{sec:perceptron} is a purely linear classifier: its decision
boundary is always a hyperplane, which limits it to solving only
\emph{linearly separable} problems.
This critical shortcoming motivated the introduction of \emph{activation
functions}---non-linear transformations applied to the weighted sum before
producing the neuron's output.

Formally, a modern neuron computes:
\begin{equation}
\label{eq:neuron-activation}
a = f\!\bigl(\mathbf{w} \cdot \mathbf{x} + b\bigr),
\end{equation}
where $f(\cdot)$ is a non-linear activation function and $a$ is the neuron's
\emph{activation}, which is subsequently transmitted to other neurons in the
network. 
One of the earliest replacements for the perceptron's step function was the
\emph{sigmoid} function:
\begin{equation}
\label{eq:sigmoid}
\sigma(z) = \frac{1}{1 + e^{-z}},
\end{equation}
which maps any real-valued input to the interval $(0,\,1)$ and is smooth
everywhere, enabling gradient-based learning algorithms such as
backpropagation~\cite{backprop}.

In modern practice, the \emph{Rectified Linear Unit}
(ReLU)~\cite{relu_glorot} is the most widely used activation function:
\begin{equation}
\label{eq:relu}
\text{ReLU}(z) = \max(0,\,z).
\end{equation}
Although piecewise linear, ReLU introduces non-linearity at the origin and
has proven to train faster than smooth alternatives on deep architectures
due to its non-saturating gradient for $z > 0$.

The insertion of a non-linear activation function after each neuron's
linear operation is what grants multi-layer networks their expressive power.
Without it, any composition of linear functions would itself remain
linear, and adding layers would provide no additional representational
capacity.
With non-linearity, each successive layer can build on the features
identified by the previous one to model increasingly complex patterns.


% ------------------------------------------------------------
\subsection{Deep Neural Networks}
\label{sec:dnn}
% ------------------------------------------------------------

When a neural network is constructed with a substantial number of layers
stacked sequentially, it is referred to as a \textbf{Deep Neural Network}
(DNN).
This ``depth'' allows the network to learn a hierarchical representation
of data: early layers extract low-level features (such as edges or textures
in images), while deeper layers compose them into higher-level abstractions
(such as objects or semantic concepts).

Computationally, each neuron in a DNN performs a weighted summation of its
inputs---an operation that primarily consists of \emph{multiply-accumulate}
(MAC) computations---followed by the application of an activation function.
The output activation is then transmitted to connected neurons in the next
layer, often in a multicast manner.
In practice, DNNs are composed of various types of specialized
layers, each implementing a specific mathematical function or
\emph{computational kernel}; due to the sheer number of
MAC operations involved and the volume of data they process,
their efficient execution has become a first-class engineering challenge,
as will be discussed in Section~\ref{sec:bg-hw-accel}.

\paragraph{Scope of this work.}
This thesis focuses exclusively on DNN \emph{inference}---the forward
evaluation of a trained model on new inputs---rather than on training, 
which involves the additional computation of gradients and parameter
updates through backpropagation.
While training workloads impose their own distinct hardware demands,
the computational kernels executed during inference (predominantly
convolutions and matrix multiplications) are the same as those in the
forward pass of training, making the analytical framework developed in
this work broadly applicable.
